\documentclass[11pt]{article}
\usepackage[margin=1in]{geometry}
\usepackage{amsmath,amssymb}
\usepackage[hidelinks]{hyperref}

\title{Disproof of Conjecture 00000001052\\
\large A maximal genus-one curve over \(\mathbb F_9\)}
\author{}
\date{}

\begin{document}
\maketitle

Use the square-field parameter \(q=3\), so maximality is over
\(\mathbb F_{q^2}=\mathbb F_9\). The conjecture asserts that every
maximal curve of genus \((q-1)^2/4=1\) is isomorphic to its specified
Hermitian curve \(H_3\). The smooth projective cubic
\[
  E:\quad Y^2Z=X^3-XZ^2
\]
is a counterexample.

\section*{Smoothness and genera}

In characteristic three, the \(X\)-derivative of
\(Y^2Z-X^3+XZ^2\) is \(Z^2\), so \(E\) is smooth on \(Z\ne0\).
At \(Z=0\) its only point is \([0:1:0]\), where the \(Z\)-derivative
is \(Y^2\ne0\). These calculations hold over the algebraic closure.
For \(H_3\), the homogeneous equation is
\[
  Y^3Z+YZ^3-X^4=0 .
\]
Its gradient in characteristic three is \((-X^3,Z^3,Y^3)\), which
cannot vanish at a projective point. Both curves are geometrically
integral. Distinct components of a reducible plane curve would meet
over the algebraic closure and give a singular point, and a repeated
component would also give singular points. The smooth plane curve
genus formula \cite[\S\,53.9]{stacks} now gives \(g(E)=1\) and
\(g(H_3)=3\).

\section*{Point counts and maximality}

Write \(k=\mathbb F_3[i]\) with \(i^2=-1\). This is a field of
nine elements because \(T^2+1\) has no root in \(\mathbb F_3\).
For \(x=a+bi\) with \(a,b\in\mathbb F_3\), one has \(x^3=a-bi\),
and hence \(x^3-x=bi\). For \(b=0\) the only solution of \(y^2=bi\)
is \(y=0\). For \(b\ne0\), the two possible right-hand sides are
\(i\) and \(-i\), both nonzero squares since
\((1-i)^2=i\) and \((1+i)^2=-i\). Each then has exactly two square
roots. The affine point count on \(E\) is \(3+6\cdot2=15\), and
the point at infinity gives \(\#E(k)=16\). This equals the
Hasse--Weil upper bound \(q^2+1+2qg(E)\), proving maximality.

For \(H_3\), \(x^4=x x^3=a^2+b^2\in\mathbb F_3\). The map
\(y\mapsto y^3+y\) is twice the real coordinate and has three
elements over every value in \(\mathbb F_3\). Thus \(H_3\) has
\(9\cdot3=27\) affine points and the same single point \([0:1:0]\)
at infinity, giving \(\#H_3(k)=28\). The two curves have different
genera and different \(k\)-point counts. In particular
\(E\not\cong H_3\), contrary to the conjecture.

The exponent in the known classification at genus \((q-1)^2/4\)
is \((q+1)/2\) \cite{fgt}. In this example \(x=iv\) and
\(y=(1+i)u\) identify the affine equation for \(E\) with
\(v^3+v=u^2\).

\section*{Lean verification}

The Lean~4.19.0 project implements arithmetic in \(k\) and proves
that its enumeration of the normalized projective charts
\([x:y:1]\), \([1:y:0]\), and \([0:1:0]\) is complete and has no
duplicate entries. It checks the two point counts and the numerical
Hasse--Weil equality. The geometric hypotheses for using that bound
and the genus formula are checked above. Run \verb|lake build| in
\verb|lean/| to verify the certificate.

\begin{thebibliography}{2}
\bibitem{stacks}
The Stacks Project, \emph{Plane curves},
\href{https://stacks.math.columbia.edu/tag/0BYA}{Tag 0BYA}.
\bibitem{fgt}
R.~Fuhrmann, A.~Garcia, and F.~Torres, \emph{On maximal curves},
\href{https://arxiv.org/abs/alg-geom/9610023}{arXiv:alg-geom/9610023}.
\end{thebibliography}

\end{document}
